\documentclass[10pt,letterpaper]{article}
\usepackage[margin=0.78in]{geometry}
\usepackage{amsmath,amssymb,booktabs,enumitem}
\usepackage[T1]{fontenc}
\usepackage{lmodern,fancyhdr}
\setlist[enumerate]{leftmargin=*,itemsep=0.32em,topsep=0.2em}
\pagestyle{fancy}\fancyhf{}
\lhead{SYDE 556/750 \quad PN1 practice test}
\rhead{30 minutes \quad 30 marks}
\cfoot{\thepage}
\setlength{\headheight}{14pt}
\newcommand{\vct}[1]{\boldsymbol{#1}}
\begin{document}
\begin{center}
{\Large\bfseries PN1 Practice Test}\\[2pt]
Neurons and Population Representation
\end{center}
\noindent\textbf{Time: 30 minutes. Total: 30 marks.} Work by hand on separate paper and show your reasoning. A basic calculator may be used. This is a practice test, not a course-issued test.\par
\noindent\textbf{Conventions.} The scalar current is \(J=\alpha e x+J^{\mathrm{bias}}\), where \(e\in\{-1,+1\}\). As in PN1, the named intercept \(\xi\) is the threshold in projected input \(e x\): firing begins when \(e x=\xi\). Maximum rate occurs at \(e x=1\). In activity matrices, rows are neurons and columns are input samples.\par

\section*{1. Scalar ReLU encoding \hfill [5 marks]}
A neuron has \(G(J)=\max(0,J)\), \(a^{\max}=120\,\mathrm{Hz}\), \(\xi=-\tfrac12\), and \(e=-1\).
\begin{enumerate}[label=\textbf{(\alph*)}]
\item Find \(\alpha\) and \(J^{\mathrm{bias}}\), showing the two current conditions you used. [2]
\item At what \(x\)-coordinate does it start firing? Explain the sign. [1]
\item Give its rates at \(x=-1,0,\tfrac12,1\), and sketch the tuning curve on \([-1,1]\). [2]
\end{enumerate}

\section*{2. Identity decoding and noise \hfill [7 marks]}
Three scalar targets and a two-neuron activity matrix are
\[
\vct{x}=\begin{bmatrix}-1\\0\\1\end{bmatrix},\qquad
A=\begin{bmatrix}0&1&1\\1&1&0\end{bmatrix},\qquad
A_{\mathrm{noisy}}=\begin{bmatrix}0&1&\tfrac32\\\tfrac32&1&0\end{bmatrix}.
\]
Use the clean \(A\) to fit decoders. For the noise-aware fit, the specified regularization amount is \(m\sigma^2=1\), where \(m=3\). Let \(\mathrm{RMSE}=\sqrt{\frac1m\sum_k(x_k-\hat x_k)^2}\).
\begin{enumerate}[label=\textbf{(\alph*)}]
\item Compute \(AA^T\) and \(A\vct{x}\), then solve \((AA^T)\vct{d}=A\vct{x}\). [2]
\item Solve \((AA^T+I)\vct{d}_{\mathrm{reg}}=A\vct{x}\). [1]
\item Fill in all four RMSEs for \(\hat{\vct{x}}^T=\vct{d}^T A_{\mathrm{test}}\). You may leave square roots unevaluated. Briefly state the tradeoff. [4]
\begin{center}\begin{tabular}{lcc}\toprule Decoder & Clean \(A\) & Noisy \(A_{\mathrm{noisy}}\)\\\midrule
Unregularized & & \\
Noise aware & & \\
\bottomrule\end{tabular}\end{center}
\end{enumerate}

\section*{3. Sources of error \hfill [4 marks]}
Illustrative mean-squared errors for different population sizes are:
\begin{center}\begin{tabular}{rccc}\toprule \(n\) & 8 & 16 & 32\\\midrule
Distortion error & 0.040 & 0.010 & 0.0025\\
Independent-noise error & 0.020 & 0.010 & 0.0050\\\bottomrule\end{tabular}\end{center}
\begin{enumerate}[label=\textbf{(\alph*)}]
\item Which row scales roughly as \(1/n\), and which as \(1/n^2\)? Explain using the effect of doubling \(n\). [2]
\item If noise standard deviation is divided by 10 with decoders fixed, what happens to noise mean-squared error? Why can independent noise error fall as more neurons are added? [2]
\end{enumerate}

\newpage
\section*{4. LIF rates and decoding \hfill [5 marks]}
For \(J>1\), \(G(J)=[\tau_{\mathrm{ref}}-\tau_{\mathrm{RC}}\ln(1-1/J)]^{-1}\); otherwise \(G(J)=0\). Use \(\tau_{\mathrm{ref}}=0.002\,\mathrm{s}\), \(\tau_{\mathrm{RC}}=0.020\,\mathrm{s}\), \(a^{\max}=100\,\mathrm{Hz}\), \(\xi=0\), \(e=+1\), and \(e^{-0.4}\approx0.6703\).
\begin{enumerate}[label=\textbf{(\alph*)}]
\item Rearrange the rate equation to express \(J_{\max}\) in terms of \(a^{\max}\), then find \(J_{\max},\alpha,J^{\mathrm{bias}}\). Use threshold current \(J_{\mathrm{th}}=1\). [3]
\item For what \(x\) is this neuron silent? How does the active part of its tuning curve differ from a ReLU line? [1]
\item For any LIF activity matrix \(B\) with \(m\) sample columns and targets \(\vct{x}\), write the noise-aware decoder normal equation. [1]
\end{enumerate}

\section*{5. Two-dimensional tuning \hfill [4 marks]}
A LIF neuron represents \(\vct{x}=(x_1,x_2)^T\) with
\(J=1+\alpha\langle\vct{e},\vct{x}\rangle\), \(\alpha>0\), and \(\vct{e}=(1/\sqrt2,-1/\sqrt2)^T\). It is silent for \(J\leq1\).
\begin{enumerate}[label=\textbf{(\alph*)}]
\item On the unit disk, give the threshold line and the silent half-disk. State the preferred direction. [2]
\item For \(\vct{x}(\theta)=(\cos\theta,\sin\theta)^T\), write \(\langle\vct{e},\vct{x}(\theta)\rangle\) as one cosine. Why is the firing-rate curve only approximately cosine-shaped? [2]
\end{enumerate}

\section*{6. Vector populations \hfill [5 marks]}
One hundred neurons represent 2D inputs. Let \(A\) contain their rates at \(m=1600\) input samples and \(X\) contain the input vectors in columns.
\begin{enumerate}[label=\textbf{(\alph*)}]
\item Give the shapes of \(A,X,D\), where \(DA\approx X\). Explain the shape of \(D\) using a \emph{single} sample of 100 activities. [2]
\item Write the noise-aware matrix normal equation for \(D\) with independent activity noise of standard deviation \(\sigma\). [1]
\item For a small illustrative population, let encoder columns be \((1,0)^T,(0,1)^T,(-1,0)^T,(0,-1)^T\) and let
\[
D=\begin{bmatrix}1&-\tfrac14&-1&\tfrac14\\-\tfrac14&1&\tfrac14&-1\end{bmatrix},\qquad
\vct{a}=\begin{bmatrix}1\\\tfrac14\\0\\0\end{bmatrix},\qquad
\vct{x}_{\mathrm{target}}=\begin{bmatrix}1\\0\end{bmatrix}.
\]
Compute the vector decoded with \(D\) and the vector decoded using the encoders as decoders. Which better preserves direction after normalization, and which better preserves magnitude in this example? Why can a broad, uniform set of encoders give a reasonable direction estimate? [2]
\end{enumerate}
\end{document}
